% from ega2-1-fr.tex
\oldpage[II]{59}

\begin{proposition}[3.4.1]
\label{II.3.4.1-fr}
Soient $Y$ un préschéma, $\mathcal{S}$ une $\mathcal{O}_Y$-Algèbre
quasi-cohérente engendrée par $\mathcal{S}_1$
(\hyperref[II.3.1.9-fr]{3.1.9})~; on suppose en outre $\mathcal{S}_1$ de type
fini. Alors $X=\operatorname{Proj}(\mathcal{S})$ est de type fini au-dessus de
$Y$.
\end{proposition}

\begin{proof}
En effet, la question étant locale sur $Y$, on peut supposer $Y$ affine
d'anneau $A$~; on a alors $\mathcal{S}=\widetilde{S}$, où
$S=\Gamma(Y,\mathcal{S})$, et par hypothèse $S$ est une $A$-algèbre engendrée
par $S_1=\Gamma(Y,\mathcal{S}_1)$, où on peut en outre supposer que $S_1$ est
un $A$-module de type fini
(\hyperref[I.1.3.9-fr]{I, 1.3.9} et
\hyperref[I.1.3.12-fr]{I, 1.3.12}). Par suite $S$ est une $A$-algèbre graduée
de type fini, et on est ramené à
(\hyperref[II.2.7.1-fr]{2.7.1, (ii)}).
\end{proof}

\begin{env}[3.4.2]
\label{II.3.4.2-fr}
Soit $\mathcal{S}$ une $\mathcal{O}_Y$-Algèbre graduée quasi-cohérente~; pour
un $\mathcal{S}$-Module gradué quasi-cohérent $\mathcal{M}$, nous considérerons
les conditions de finitude suivantes~:
\begin{enumerate}
  \item[(\textbf{TF})] \emph{Il existe un entier $n$ tel que le
    $\mathcal{S}$-Module $\bigoplus_{k\geq n}\mathcal{M}_k$ soit de type fini.}
  \item[(\textbf{TN})] \emph{Il existe un entier $n$ tel que
    $\mathcal{M}_k=0$ pour $k\geq n$.}
\end{enumerate}
Si $\mathcal{M}$ vérifie (\textbf{TN}), on a
$\widetilde{\mathcal{M}}=0$, la question étant locale sur $Y$
(\hyperref[II.2.7.2-fr]{2.7.2}).

Soient $\mathcal{M}$, $\mathcal{N}$ deux $\mathcal{S}$-Modules gradués
quasi-cohérents~; on dit qu'un homomorphisme
$u:\mathcal{M}\to\mathcal{N}$ de degré $0$ est
(\textbf{TN})-\emph{injectif} (resp. (\textbf{TN})-\emph{surjectif},
(\textbf{TN})-\emph{bijectif}) s'il existe un entier $n$ tel que
$u_k:\mathcal{M}_k\to\mathcal{N}_k$ soit injectif (resp. surjectif, bijectif)
pour $k\geq n$~; alors
$\widetilde{u}:\widetilde{\mathcal{M}}\to\widetilde{\mathcal{N}}$ est injectif
(resp. surjectif, bijectif) en vertu de
(\hyperref[II.2.7.2-fr]{2.7.2}), la question étant locale sur $Y$, et compte
tenu de (\hyperref[I.1.3.9-fr]{I, 1.3.9})~; lorsque $u$ est
(\textbf{TN})-bijectif, on dit aussi que $u$ est un
(\textbf{TN})-\emph{isomorphisme}.
\end{env}

\begin{proposition}[3.4.3]
\label{II.3.4.3-fr}
Soient $Y$ un préschéma, $\mathcal{S}$ une $\mathcal{O}_Y$-Algèbre graduée
quasi-cohérente engendrée par $\mathcal{S}_1$, $\mathcal{S}_1$ étant supposé
de type fini. Soit $\mathcal{M}$ un $\mathcal{S}$-Module gradué
quasi-cohérent.
\begin{enumerate}
  \item[\rm(i)] Si $\mathcal{M}$ vérifie la condition (\textbf{TF}),
    $\widetilde{\mathcal{M}}$ est de type fini.
  \item[\rm(ii)] Supposons que $\mathcal{M}$ vérifie (\textbf{TF})~; pour que
    $\widetilde{\mathcal{M}}=0$, il faut et il suffit que $\mathcal{M}$
    vérifie (\textbf{TN}).
\end{enumerate}
\end{proposition}

\begin{proof}
Les questions étant locales sur $Y$, on est ramené au cas où $Y$ est affine
d'anneau $A$, $\mathcal{S}=\widetilde{S}$, où $S$ est une $A$-algèbre graduée
telle que l'idéal $S_+$ soit de type fini,
$\mathcal{M}=\widetilde{M}$ où $M$ est un $S$-module gradué~; la proposition
résulte alors de (\hyperref[II.2.7.3-fr]{2.7.3}).
\end{proof}

\begin{theorem}[3.4.4]
\label{II.3.4.4-fr}
Soient $Y$ un préschéma, $\mathcal{S}$ une $\mathcal{O}_Y$-Algèbre graduée
quasi-cohérente engendrée par $\mathcal{S}_1$, $\mathcal{S}_1$ étant supposé
de type fini~; soit $X=\operatorname{Proj}(\mathcal{S})$. Pour tout
$\mathcal{O}_X$-Module quasi-cohérent $\mathcal{F}$, l'homomorphisme canonique
$\beta$ (\hyperref[II.3.3.4-fr]{3.3.4}) est un isomorphisme.
\end{theorem}

\begin{proof}
Notons d'abord que $\beta$ est défini en vertu de
(\hyperref[II.3.4.1-fr]{3.4.1}). Pour voir que $\beta$ est un isomorphisme,
on est ramené au cas où $Y$ est affine d'anneau $A$,
$\mathcal{S}=\widetilde{S}$, où $S$ est une $A$-algèbre graduée engendrée par
$S_1$ et $S_1$ est un $A$-module de type fini. Il suffit alors d'appliquer
(\hyperref[II.2.7.5-fr]{2.7.5}).
\end{proof}

\begin{corollary}[3.4.5]
\label{II.3.4.5-fr}
Sous les hypothèses de (\hyperref[II.3.4.4-fr]{3.4.4}), tout
$\mathcal{O}_X$-Module quasi-cohérent $\mathcal{F}$ est isomorphe à un
$\mathcal{O}_X$-Module de la forme $\widetilde{\mathcal{M}}$, où
$\mathcal{M}$ est un $\mathcal{S}$-Module gradué quasi-cohérent. Si en outre
$\mathcal{F}$ est de type fini, et si on suppose que $Y$ est un schéma
quasi-compact, ou que l'espace sous-jacent à $Y$ est noethérien, on peut
supposer $\mathcal{M}$ de type fini.
\end{corollary}

